\subsection{Round two results}
\label{sec:r2-results}

\begin{table}[t]
\caption{Round two: geomean ratios vs the unseeded baseline at scale~1 for the soft arms (Urd from
round one for reference). The last column is the number of designs where the routed wirelength is
below the baseline. WL is routed wirelength, Chan.\ the minimum channel width, Time the
placement runtime, $T_0$ the anneal's start temperature that VPR derives from the
starting placement. Below 1.0 is better (except $T_0$, a diagnostic).}
\label{tab:r2qor}
\centering
\footnotesize\setlength{\tabcolsep}{3pt}
\begin{tabular}{lcccccc}
\toprule
Seed & WL & Delay & Chan.\ & Time & $T_0$ & Better \\
\midrule
Urd (round 1) & 1.026 & 1.047 & 0.933 & 1.061 & 2.96 & 2/9 \\
Words & 1.013 & 0.997 & 0.945 & 1.257 & 5.83 & 3/9 \\
Words, inter only & 1.018 & 1.006 & 0.944 & 1.272 & 7.04 & 3/9 \\
Slices & 0.969 & 0.998 & 1.002 & 1.278 & 8.92 & 7/9 \\
Hier & 1.010 & 1.006 & 0.968 & 1.257 & 4.40 & 4/9 \\
Words, partial & 0.985 & 1.008 & 0.936 & 1.262 & 5.42 & 5/9 \\
Random start & 0.975 & 1.003 & 0.996 & 1.274 & 6.94 & 5/9 \\
Urd, permuted & 1.035 & 1.011 & 0.941 & 1.241 & 2.91 & 2/9 \\
\bottomrule
\end{tabular}
\end{table}

\textbf{RQ-N7, bitwise seeds (9 designs, one VPR seed each).}
H7a is inconclusive: \code{words} over Urd is 0.986 in wirelength, better on 5 of
9 designs (sign test $p=1.000$).
H7b is inconclusive: \code{slices} over \code{words} is 0.957 in wirelength
(better on 7 of 9, $p=0.180$) and 1.060 in channel width
(better on 3 of 9).
H7c is not supported: \code{hier} ranks 8 of 11 families at
1.010; the lowest is \code{slices} at 0.969.
H7d is supported: \code{words\_partial} over \code{words} is 0.973
(better on 6 of 9, $p=0.508$).
Ties (identical results) are left out of the sign tests: 0 for H7a, 0 for H7b and
0 for H7d. A family with nothing to rearrange reproduces its reference exactly (a design whose words are
all single-CLB, or that has one module): \code{slices} equals \code{words} on 0
designs, \code{hier} equals Urd on 0 and \code{words\_partial} equals \code{words} on
0.
The best new family by geomean wirelength is \code{slices}, so the coarse-random, hard and
crossover arms below use it.

\textbf{RQ-N4, clustering versus coarse layout.} H4 is not supported. For Urd the wirelength
ratios are random start 0.975, permuted cluster positions 1.035, quotient
layout 1.026; for \code{slices} they are 0.975, n/a and 0.969. The analytic baseline is 1 by definition, so a random start above 1 means VPR's analytic
start is better than a random one, and the gap between the permuted and the quotient arm is what
the coarse layout is worth beyond the clustering itself.

\textbf{RQ-N1, does VPR rediscover the hierarchy?} H1 is inconclusive. Table~\ref{tab:r2n1}
gives medians over designs: the final placement is above the random null in NMI on 8 of 13
designs for modules and 9 of 15 for words, and in compactness on 10
and 14 (more than two null standard deviations). Final and analytic placements are close, and the final
NMI is at least the analytic one on 5 of 13 (modules) and 7 of 15
(words) designs. With many small labels the NMI of a random placement is itself high (word null
0.66), so compactness is the cleaner signal.

\begin{table}[t]
\caption{RQ-N1: medians over designs of the grid-partition NMI and label compactness of the
analytic start, VPR's final placement and a random null (20 permutations). Designs with a
single module have no module label and are left out of that row.}
\label{tab:r2n1}
\centering
\footnotesize\setlength{\tabcolsep}{3pt}
\begin{tabular}{lccc}
\toprule
Labels & $n$ & NMI (start/final/null) & Comp.\ (start/final/null) \\
\midrule
Module & 13 & 0.25/0.23/0.10 & 0.28/0.26/0.10 \\
Word & 15 & 0.72/0.73/0.66 & 0.33/0.32/0.08 \\
\bottomrule
\end{tabular}
\end{table}

\textbf{RQ-N2, does structure predict the gain?} H2 is not supported: 0 of
9 correlations reach $p<0.05$ (no multiple-comparison correction; $n=15$ designs, so only a
large $\rho$ can be told from chance). Table~\ref{tab:r2n2} lists Spearman's $\rho$ (and $p$).

\begin{table}[t]
\caption{RQ-N2: Spearman $\rho$ ($p$) between per-design structure and seed gain
(1 - wirelength ratio): Urd round one, best of seven round-one families, best new family. The
inverse Rent exponent is $-p$.}
\label{tab:r2n2}
\centering
\footnotesize\setlength{\tabcolsep}{3pt}
\begin{tabular}{lccc}
\toprule
Gain & $Q$ & $z(Q)$ & $-p_{\text{Rent}}$ \\
\midrule
Urd, round one & -0.17 (0.533) & -0.45 (0.089) & 0.12 (0.657) \\
Best round one & -0.09 (0.742) & -0.24 (0.390) & -0.16 (0.558) \\
Best new family & 0.15 (0.700) & -0.37 (0.332) & 0.43 (0.244) \\
\bottomrule
\end{tabular}
\end{table}
